
% written in plaintex
\input header
\def\PSL{{\rm PSL}}
\def\PD{\rho}
\def\SD{\sigma}
\footline={\hss\tenrm i-\folio\hss}\pageno=1

%table of content format
\def\leaderfill{{\leaders\hbox to 1em{\hss.\hss}\hfill}\hfil}
\def\eleaderfill{\hfill\hbox to 1em{\hss\hss}\hfill}

\vbox{}\vskip -2\baselineskip

\centerline{\bf DYNAMICS IN ONE COMPLEX VARIABLE}\smallskip
\centerline{\bf Introductory Lectures}\smallskip
\centerline{(Partially revised version of 9-5-91)}
\smallskip

\centerline{\bf John Milnor}\smallskip
% (SUNY Stony Brook and IAS Princeton)}\smallskip % \medskip

\centerline{\bit Institute for Mathematical Sciences, SUNY, Stony Brook NY}
\medskip

\centerline{\bf Table of Contents.}\smallskip
\line{Preface\leaderfill }
\line{Chronological Table\leaderfill }
\smallskip
\line{{\bf Riemann Surfaces}.\eleaderfill}
\line{\kern15pt 1.  Simply Connected Surfaces\leaderfill (9 pages)}
\line{\kern15pt 2. The Universal Covering, Montel's Theorem\leaderfill (5)}
\smallskip
\line{{\bf The Julia Set}.\eleaderfill}
\line{\kern15pt 3. Fatou and Julia: Dynamics
on the Riemann Sphere\leaderfill (9)}
\line{\kern15pt 4. Dynamics on Other Riemann Surfaces\leaderfill (6)}
\line{\kern15pt 5. Smooth Julia Sets\leaderfill (4)}
\smallskip
\line{{\bf Local Fixed Point Theory}.\eleaderfill}
\line{\kern15pt 6. Attracting and Repelling Fixed Points\leaderfill (7)}
\line{\kern15pt 7. Parabolic Fixed Points: the Leau-Fatou Flower\leaderfill (8)}
\line{\kern15pt 8. Cremer Points and Siegel Disks \leaderfill (11)}
\smallskip
\line{{\bf Global Fixed Point Theory}.\eleaderfill}
\line{\kern15pt 9. The Holomorphic Fixed Point Formula\leaderfill (3)}
\line{\kern10pt 10. Most Periodic Orbits Repel\leaderfill (3)}
\line{\kern10pt 11. Repelling Cycles are Dense in \[J\]\leaderfill (5)}
\smallskip
\line{{\bf Structure of the Fatou Set}.\eleaderfill}
\line{\kern10pt 12. Herman Rings\leaderfill (4)}
\line{\kern10pt 13. The
Sullivan Classification of Fatou Components\leaderfill (4)}
\line{\kern10pt 14. Sub-hyperbolic and hyperbolic Maps\leaderfill (7)}
\smallskip
\line{{\bf Carath\'eodory Theory}.\eleaderfill}
\line{\kern10pt 15. Prime Ends\leaderfill (5)}
\line{\kern10pt 16. Local Connectivity\leaderfill (4)}
\smallskip
\line{{\bf Polynomial Maps}.\eleaderfill}
\line{\kern10pt 17. The Filled Julia Set \[K\,$\leaderfill (4)}
\line{\kern10pt 18. External Rays and Periodic Points\leaderfill (8)}
\smallskip
\line{Appendix A. Theorems from Classical Analysis\leaderfill (4)}
\line{Appendix B. Length-Area-Modulus Inequalities\leaderfill (6)}
\line{Appendix C. Continued Fractions\leaderfill (5)}
\line{Appendix D. Remarks on Two Complex Variables\leaderfill (2)}
\line{Appendix E. Branched Coverings and Orbifolds\leaderfill (4)}
\line{Appendix F. Parameter Space\leaderfill (3)}
\line{Appendix G. Remarks on Computer Graphics\leaderfill (2)}
\smallskip
\line{References\leaderfill (7)}
\line{Index\leaderfill (3)}



\eject

\centerline{\bf PREFACE.}\bigskip

These notes will study the dynamics of iterated holomorphic
mappings from a\break Riemann surface to itself, concentrating on
the classical case of rational maps of the\break Riemann sphere.
They are based on introductory lectures
given at Stony Brook during the Fall Term of 1989-90. I am grateful
to the audience
for a great deal of constructive criticism, and
to Branner, Douady, Hubbard, and Shishikura who
taught me most of what I know in this field. The surveys by
Blanchard, Devaney, Douady, Keen, and Lyubich have been extremely useful, and
are highly recommended
to the reader. (Compare the list of references at the end.) Also, I want to
thank A. Poirier for his criticisms of my first draft.\smallskip

This subject is large and rapidly growing. These lectures are intended to
introduce the reader to some key ideas in
the field, and to form a basis for further study.
The reader is assumed to be familiar with the rudiments of complex variable
theory and of two-dimensional differential geometry. % Compare the list
\bigskip
\rightline{John Milnor, Stony Brook, February 1990}\medskip
\vfil\eject  % \bigskip

\centerline{\bf CHRONOLOGICAL TABLE}\bigskip

Following is a list of some of the founders of the field of complex dynamics.
\smallskip

$$\vbox{\settabs 2 \columns
\+Hermann Amandus Schwarz & 1843--1921\cr
\+Henri Poincar\'e & 1854--1912\cr
\+Gabriel Ko$\!$enigs	& 1858--1931\cr
\+L\'eopold Leau &	1868--1940$\pm$\cr
% \+L. \`E. B\"otk$\!$her  & ??\cr %(= L. \`E. B\"otk$\!$her = &\cr
\+Lucyan Emil B\"ottcher  & 1872-- ~?\cr
\+Constantin Carath\'eodory & 1873--1950\cr
\+Samuel Latt\`es & 1875$\pm$-1918\cr
\+Paul Montel   & 1876--1975\cr
\+Pierre Fatou &	1878--1929\cr
\+Paul Koebe & 1882--1945\cr
\+Gaston Julia &        1893--1978 \cr
\+Carl Ludwig Siegel &  1896--1981\cr
\+Hubert Cremer &	1897--1983\cr
\+Charles Morrey & 1907--1984\cr}$$
\smallskip

\noindent Among the many present day workers in the field, let me mention
a few whose work is emphasized in these notes:
Lars Ahlfors (1907),
Lipman Bers (1914),
Adrien Douady (1935),
Vladimir I. Arnold (1937),
Dennis P. Sullivan (1941),
Michael R. Herman (1942),
Bodil Branner (1943),
John Hamal Hubbard    (1945),
William P. Thurston (1946),
Mary Rees (1953),
Jean-Christophe Yoccoz (1955),
Mikhail Y. Lyubich (1959),
and Mitsuhiro Shishikura (1960).
% \bigskip
\vfil
%\noindent ------------

%\noindent * B\"ottcher was born in
%Warsaw, took his doctorate in Leipzig (1898), and was active in Lvov
%and Kazan. The Russian form of his name is ~{\cyr L. \`E. B\"etkher\cdprime}.

\eject
\footline={\hss\tenrm 1-\folio\hss}\pageno=1


\centerline{\bf RIEMANN SURFACES}\medskip


\centerline{\bf \S1. Simply Connected Surfaces.}\medskip

The first two sections will present an overview of well known material.
\smallskip

By a {\bit Riemann surface} we mean a connected complex analytic manifold
of complex dimension one. Two such surfaces \[S\] and \[S'\] are
{\bit conformally isomorphic} if there is a homeomorphism
from \[S\] onto \[S'\] which is holomorphic, with holomorphic inverse.
(Thus our conformal maps must always preserve orientation.) According
to \Poincare and Koebe, there are only three simply connected Riemann surfaces,
up to isomorphism.

{\QP{\bf 1.1. Uniformization Theorem.} \it Any simply connected Riemann
surface is conformally isomorphic either

\noindent $(a)$ to the plane \[\C\]
consisting of all complex numbers \[z=x+iy\],

\noindent $(b)$ to the open unit disk
\[D\subset\C\] consisting of all \[z\] with \[|z|^2=x^2+y^2<1\], or

\noindent $(c)$ to the
Riemann sphere \[\hat\C\] consisting of \[\C\] together with a point at
infinity.\medskip}

The proof, which is quite difficult, may be found in Springer, or Farkas \&
Kra, or Ahlfors [A2], or in Beardon.\QED

For the rest of this section, we will discuss these three surfaces in more
detail. We begin with a study of the unit disk \[D\].

{\QP{\bf 1.2. Schwarz Lemma.} {\it If \[f:D\to D\] is a holomorphic map with
\[f(0)=0\], then the derivative at the origin satisfies \[|f'(0)|\le 1\].
If equality holds, then \[f\] is a rotation about the origin, that is
\[f(z)=\lambda z\] for some constant \[\lambda
=f'(0)\] on the unit circle. In particular,
it follows that \[f\] is a conformal automorphism of \[D\]. On the other
hand, if \[|f'(0)|<1\], then \[|f(z)|<|z|\] for all \[z\ne 0\], and
\[f\] is not a conformal automorphism.}\medskip}

{\bf Proof.} We use the {\bit maximum modulus principle}, which
asserts that a non-constant
holomorphic function cannot attain its maximum absolute value at 
any interior point of its region of definition. First note that the
quotient function \[g(z)=f(z)/z\] is well defined and holomorphic
throughout the disk \[D\]. Since \[|g(z)|<1/r\] when \[|z|=r<1\], it follows
by the maximum modulus principle that \[|g(z)|< 1/r\] for all \[z\] in the
disk \[|z|\le r\]. Taking the limit as \[r\to 1\], we see that \[|g(z)|\le
1\] for all \[z\in D\]. Again by the maximum modulus principle, we see that
the case \[|g(z)|=1\], with \[z\] in the open disk, can occur only if the
function \[g(z)\] is constant. If we exclude this case \[f(z)/z=c\],
then it follows that \[|g(z)|=|f(z)/z|<1\] for all \[z\ne 0\], and similarly
that \[|g(0)|=|f'(0)|<1\]. Evidently the composition of two such maps
must also satisfy \[|f_1(f_2(z))|<|z|\], and hence cannot be the identity
map of \[D\]. \QED

{\bf 1.3. Remarks.} The Schwarz Lemma was first proved in this generality
by Carath\'eo\-dory. If we map the disk \[D_r\] of radius \[r\] into the disk
\[D_s\] of radius \[s\],
with \[f(0)=0\], then a similar argument shows that \[|f'(0)|
\le s/r\]. Even if we drop the condition that \[f(0)=0\], we certainly get the
inequality
$$      |f'(0)|\le 2s/r\qquad{\rm whenever}\qquad f(D_r)\subset D_s     $$
since the difference \[f(z)-f(0)\] takes values in \[D_{2s}\].
(In fact the extra factor of 2 is unnecessary. Compare Problem 1-2.)
One easy corollary
is the {\bit Theorem of Liouville}, which says that a bounded function
which is defined and holomorphic everywhere on \[\C\] must be constant.
Another closely related statement is the following.

{\QP{\bf 1.4. Theorem of Weierstrass.} \it If a sequence of holomorphic
functions converges uniformly, then their derivatives also converge uniformly,
and the limit function is itself holomorphic.\smallskip}

\noindent In fact the convergence of first derivatives follows easily from the
discussion above. For the proof of the full theorem, see for example [A1].
\smallskip

It follows from this discussion that our three model surfaces
really are distinct. For there are natural inclusion maps \[D\to\C\to\hat\C\];
yet it follows from the maximum modulus principle and Liouville's Theorem
that every holomorphic map \[\hat\C\to\C\] or %\break
\[\C\to D\] must be constant.\medskip

It is often more convenient to work with the {\bit upper half-plane}
\[H\], consisting of all complex numbers \[w=u+iv\] with \[v>0\].

{\QP{\bf 1.5. Lemma.} {\it The half-plane \[H\] is conformally isomorphic to
the disk \[D\] under the holomorphic mapping \[z=(i-w)/(i+w)\], with
inverse \[w=i(1-z)/(1+z)\].}\medskip}

{\bf Proof.} We have \[|z|^2<1\] if and only if \[|i-w|^2=
u^2+(1-2v+v^2)\] is less than
\[|i+w|^2=u^2+(1+2v+v^2)\], or in other words if and only if \[0<v\]. \QED

{\QP{\bf 1.6. Lemma.} {\it Given any point \[z_0\] of \[D\], there exists
a conformal automorphism \[f\] of \[D\] which maps \[z_0\] to the origin.
Furthermore, \[f\] is uniquely determined up to composition with a rotation
which fixes the origin.}\medskip}

{\bf Proof.} Given any two points \[w_1\] and \[w_2\] of the upper
half-plane \[H\], it is easy to check that there
exists a automorphism of the form \[w\mapsto a+bw\] which carries \[w_1\]
to \[w_2\]. Here the coefficients \[a\] and \[b\]
are to be real, with \[b>0\]. Since \[H\]  is conformally isomorphic to
\[D\], it follows that there exists a conformal automorphism of \[D\]
carrying any given point to zero. As noted above, the Schwarz Lemma
implies that
any automorphism of \[D\] which fixes the origin is a rotation. \QED

{\QP{\bf 1.7. Theorem.} \it The group \[G(H)\] consisting of all
conformal automorphisms of the upper half-plane \[H\] can be identified
with the group of all fractional linear transformations \[w\mapsto
(aw+b)/(cw+d)\] with real coefficients and with determinant \[ad-bc>0\].
\medskip}

If we normalize so that \[ad-bc=1\], then these coefficients are
well defined up to a simultaneous change of sign. {\it
Thus \[G(H)\] is isomorphic to the group \[\PSL(2,\R)\], consisting of all
\[2\times 2\] real matrices with determinant \[+1\] modulo the subgroup
\[\{\pm I\}\].}  To prove 1.7, it is only necessary to
note that this group acts transitively on \[H\], and that it contains the
group of ``rotations''
$$	g(w)= (w\cos\theta+\sin\theta)/(-w\sin\theta+\cos\theta)\,,
 \eqno (1)$$
which fix the point \[g(i)=i\] with \[g'(i)=e^{2i\theta
}\]. By 1.6, there can be no further
automorphisms. (Compare Problem 1-2.)\QED\smallskip

Next we introduce the {\bit Poincar\'e metric} on the half-plane \[H\].

{\QP{\bf 1.8. Theorem.} {\it There exists one and, up to multiplication by
a constant, only one Riemannian metric on the half-plane \[H\] which is invariant
under every conformal automorphism of \[H\].}\medskip}

It follows immediately that the same statement is true for the disk
\[D\], or for any other Riemann surface which is conformally isomorphic
to \[H\].\smallskip

{\bf Proof of 1.8.} A Riemannian metric \[ds^2=g_{11}dx^2
+2g_{12}dxdy+g_{22}dy^2\] is said to be {\bit conformal}$\,$ if
\[g_{11}=g_{22}\] and \[g_{12}=0\], so that the matrix
\[[g_{ij}]\], evaluated at any point \[z\],
is some multiple of the identity matrix. Such a metric can be
written as\break \[ds^2=\gamma(x+iy)^2(dx^2+dy^2)\], or briefly as \[ds=\gamma(z)
|dz|\], where the function \[\gamma(z)\] is smooth and strictly positive.
By definition, such a metric is {\bit invariant} under a conformal
automorphism \[z'=f(z)\] if and only if it satisfies the identity
\[\gamma(z')|dz'|=\gamma(z)|dz|\], or in other words.
$$	\gamma(f(z))=\gamma(z)/|f'(z)|\,.	\eqno (2)	$$
Equivalently, we may say that \[f\] is an {\bit isometry} with respect to
the metric.

As an example, suppose that a conformal metric \[\gamma(w)
|dw|\] on the upper half-plane
is invariant under every linear automorphism \[f(w)=a+bw\]. Then we must
have\break \[\gamma(a+bi)=\gamma(i)/b\,\]. If we normalize by setting \[\,
\gamma(i)
=1\,\], then we are led to the\break formula \[\gamma(u+iv)=1/v\], or
in other words
$$	ds=|dw|/v\qquad {\rm for}\quad w=u+iv\in H \,. \eqno(3) $$
{\it In fact, the metric defined by this formula is invariant under every
conformal automorphism \[g\] of \[H\].} For, if we select some arbitrary point
\[w_1\in H\] and set \[g(w_1)=w_2\], then \[g\]
can be expressed as the composition of a linear automorphism of the form
\[g_1(w)= a+bw\] which maps \[w_1\] to \[w_2\]
and an automorphism \[g_2\] which fixes \[w_2\].
We have constructed the metric (3) so that \[g_1\] is an isometry,
and it follows from
1.2 and 1.5 that \[|g'_2(w_2)|=1\], so that \[g_2\] is an isometry at \[w_2\].
Thus our metric is invariant at an arbitrarily chosen point under an arbitrary
automorphism.

To complete the proof of 1.8, we must show that a metric which is invariant
under all automorphisms
of \[H\] is necessarily conformal. For this purpose, it is
convenient to pass to the equivalent problem on \[D\]. In fact
a brief computation shows that any metric on \[D\] which is 
preserved by all rotations about the origin must be conformal
at the origin. Further details will be left to the reader. \QED

{\bf Definition.} This metric \[ds=|dw|/v\] is called
the {\bit Poincar\'e metric} on the upper half-plane
\[H\]. It can be shown that this is
the unique conformal metric on \[H\] which is complete, with
constant Gaussian curvature equal to \[-1\].
(Compare Problems 1-9 and 2-2.)
The corresponding expression on \[D\] is
$$ ds = 2|dz|/(1-|z|^2)\qquad {\rm for}\quad z=x+iy\in D\,, \eqno (4)$$
as can be verified using 1.5 and (2).\smallskip

{\bf Caution:} Some authors call \[|dz|/(1-|z|^2)\] the Poincar\'e
metric on \[D\], and correspondingly call
\[{1\over 2}|dw|/v\] the Poincar\'e metric on \[H\].
These modified metrics have constant Gaussian curvature equal to \[-4\].
\smallskip

Thus there is a preferred Riemannian metric \[ds\] on \[D\] or on \[H\].
More generally, if \[S\] is any Riemann surface which is conformally
isomorphic to \[D\], then there is a corresponding Poincar\'e metric
\[ds\] on \[S\].
Hence we can define the {\bit Poincar\'e distance} \[\PD(z_1\,,\, z_2)
=\PD_S(z_1\,,\, z_2)\] between two points of \[S\]
as the minimum, over all paths \[P\] joining
\[z_1\] to \[z_2\], of the integral \[\,\int_P\,ds\]. 

{\QP{\bf 1.9. Lemma.} \it The disk \[D\]
with this \Poincare metric is {\bit complete}. That is:

\noindent $(a)$ every Cauchy sequence with respect to the metric \[\rho_D\] converges,

\noindent $(b)$ every closed neighborhood \[\;N_r(z_0\,,\,\rho_D)\;=\;\{z\in D : \rho_D
(z,z_0)\le r\}\;\] is a compact subset of \[D\], and

\noindent$(c)$ any path leading from \[z_0\] to a point of the boundary
\[\partial D=\bar D\ssm D\subset \hat\C\] has infinite \Poincare length.

\noindent Furthermore, any two points of \[D\] are joined by a unique
minimal geodesic.\medskip}

(Compare Willmore.) Equivalently, we
have exactly the same statements for the half-plane \[H\].\smallskip

{\bf Proof.} Given any two points of \[D\] we can first choose a conformal
automorphism which moves the first point to the origin and the second
to some point \[\xi\] on the positive real axis. For any path \[P\]
between \[0\] and \[\xi\] within \[D\] we have
$$	\int_P ds\;=\;\int_P\;{2|dz|\over 1-|z|^2}\;\ge\;\int_P\,{2|dx|\over
 1-x^2}	\;\ge\;\int_{[0,\xi]}\;{2dx\over 1-x^2}\;=\;\log{1+\xi\over 1-\xi}\;,
	$$
with equality if and only if \[P\] is the straight line segment \[[0,\xi]\].
For any \[z\in D\], it follows easily  that the \Poincare distance \[\rho
=\rho_D
(0,z)\] is equal to \[\log\big((1+|z|)/(1-|z|)\big)\], and that the straight line
segment from \[0\] to \[z\] is the unique minimal \Poincare geodesic. Clearly
\[\rho\to\infty\] as \[|z|\to 1\], which proves (c). The proof of (b), and
hence of (a) is now straightforward.\QED\smallskip

The \Poincare metric
has the marvelous property of never increasing under holomorphic maps.

{\QP {\bf 1.10. Theorem of Pick.} {\it If \[f:S\to T\] is a holomorphic
map between
Riemann surfaces, both of which are conformally isomorphic to \[D\], then
$$	  \PD_{T}(f(z_1), \,f(z_2))\;\le\;\PD_{S}(z_1,\,z_2) \,.	$$
Furthermore, if equality holds for some
\[z_1\ne z_2\] in \[S\], then \[f\] must be
a conformal isomorphism from \[S\] onto \[T\].}\medskip}

{\bf Proof.} Join \[z_1\] to \[z_2\] by a geodesic of length equal to
the distance \[\PD_{S}(z_1\,,
\,z_2)\]. Let \[ds\] be the element of \Poincare length along this
geodesic, and let \[ds'\] be the element of length along the image curve in
\[T\]. It follows from the Schwarz Lemma and the definition of the
\Poincare metric that \[|ds'/ds|\le 1\], with equality if and only if \[f\]
is a conformal isomorphism; and the conclusion follows.\QED

Note that this distance depends sharply on the choice of \[S\].
As an example, suppose that \[U\subset D\] is a
simply connected open subset with \[U\ne D\]. Then \[U\] is conformally
isomorphic to \[D\] by the Riemann Mapping Theorem. Applying 1.10 to the
inclusion\break \[U\to D\], we see that
$$	\PD_D(z_1\,,\,z_2)\;<\;	\PD_U(z_1\,,\,z_2)	$$
for any two distinct points \[z_1\] and \[z_2\] in \[U\].\smallskip

{\bf 1.11. Remark.} In the special case of a holomorphic map from \[S\]
to itself, if the sharper inequality
$$	\PD(f(z_1),f(z_2))\le c\PD(z_1\,,\,z_2)	$$
is satisfied for all \[z_1\] and \[z_2\], where \[0<c<1\] is constant,
then \[f\] necessarily has a unique fixed point. However,
the example \[z\mapsto z^2\] on the unit disk shows that a map with a
unique fixed point need not satisfy this sharper inequality; and
the example \[w\mapsto w+i\] on the upper half-plane shows that map which
simply reduces \Poincare distance need not have any fixed point. (See also
Problem 1-1.)\medskip

Next let us consider the {\bit Riemann sphere} \[\hat\C\],
that is the compact Riemann surface consisting of the complex
numbers \[\C\] together with a single point at infinity. The conformal
structure on this complex manifold can be specified by using
the usual coordinate \[z\] as uniformizing parameter in the finite plane
\[\C\], and using \[\zeta=1/z\] as uniformizing parameter in \[\hat\C\ssm
\{0\}\].

When studying \[\hat\C\], we need some {\bit spherical metric} \[\SD\]
which is adapted to the topology, so that the point at infinity has
finite distance from other points of \[\hat\C\]. The precise choice of such
a metric does not matter for our purposes. However, one good choice
would be the distance function \[\SD(z_1\,,\,z_2)\] which is associated with
the Riemannian metric
$$	ds=2|dz|/(1+|z|^2)\,.	\eqno (6)	$$
This metric is smooth and well behaved, even in a neighborhood
of the point at infinity, and has constant Gaussian curvature \[+1\].
It corresponds to the standard metric on the unit sphere in \[\R^3\]
under stereographic projection.
Note that the map \[z\mapsto 1/z\] is an isometry for this metric.
(However, it is certainly not true that every conformal self-map of \[\hat\C\]
is an isometry.)\smallskip

The group \[G(\hat\C)\] consisting of all conformal automorphisms of \[\hat
\C\] can be described as follows. By definition,
an automorphism \[g\in G(\hat\C)\] is called an {\bit involution}
if the composition \[g\circ g\] is the identity map of \[\hat\C\].

{\QP{\bf 1.12. Theorem.} \it 
Every conformal automorphism \[g\] of \[\hat\C\] can be
expressed as a {\bit fractional linear transformation} or {\bit M\"obius
transformation}
$$	g(z)\;=\; (az+b)/(cz+d)\,,	$$
where the coefficients are complex numbers with determinant \[ad-bc\ne 0\].
Every non-identity automorphism of \[\hat\C\] either has two distinct fixed
points or one double fixed point in \[\hat\C\]. In general,
two non-identity elements of \[G(\hat\C)\] commute if and only if they have
precisely the same fixed points: the only exceptions to this statement
are provided by pairs of commuting involutions.\medskip}

Here by a ``double'' fixed point we mean one at which the derivative \[g'(z)\]
is equal to \[+1\].
If we normalize so that \[ad-bc=1\], then the coefficients are
well defined up to a simultaneous change of sign.
{\it Thus the group \[G(\hat\C)\] of conformal
automorphisms can be identified with the group \[\PSL(2,\C)\] consisting of
all complex matrices with determinant \[+1\] modulo the subgroup \[\{\pm
I\}\].}\smallskip

{\bf 1.13. Remark.} The group \[G(\C)\] of all conformal automorphisms
of the complex plane can be identified with the subgroup of \[G(\hat\C)\]
consisting of automorphisms \[g\] which fix the point \[\infty\], since
every conformal automorphism of \[\C\] extends uniquely
to a conformal automorphism of \[\hat\C\]. (Compare Ahlfors [A1, p.124] .)
It follows easily that \[G(\C)\] consists
of all affine transformations
$$	g(z)\;=\;\lambda z+c	 $$
with complex coefficients \[\lambda\ne 0\] and \[c\].\smallskip

{\bf Proof of 1.12.} Clearly \[G(\hat\C)\] contains this group of fractional
linear transformations as a subgroup. After composing the given \[g\in G(\hat
\C)\] with a suitable element of this subgroup, we may assume that \[g(0)=0\]
and \[g(\infty)=\infty\]. But then the quotient \[g(z)/z\] is a bounded
holomorphic function, hence constant by Liouville's Theorem (\S1.3). Thus
\[g\] is linear, and hence itself is an element of \[\PSL(2,\C)\].

The fixed points of a fractional linear transformation can be
determined by solving a quadratic equation, so it is easy to check that there
must be at least one and
at most two distinct solutions in the extended plane \[\hat\C\].
In particular, if an automorphism of \[\hat\C\] 
fixes three distinct points, then it must be the identity map.

In general, if two automorphisms \[g\] and \[h\] commute, we
can say that \[g\] maps every fixed point of \[h\] to a fixed point of \[h\],
and that \[h\] maps every fixed point of \[g\] to a fixed point of \[g\].
However, this leaves open the possibility that \[g\] interchanges the two fixed
of \[h\] and that \[h\] interchanges the two fixed points of \[g\].
If this is indeed the case, then both \[g\circ g\]
and \[h\circ h\] have at least four fixed points, hence both must equal
the identity map. (An example of this phenomenon is provided by the two
commuting involutions \[g(z)= -z\] with fixed points \[\{0,\infty\}\], and
\[h(z)= 1/z\] with fixed points \[\{\pm 1\}\].) If we exclude this possibility,
then elements which commute must have exactly the same fixed points.

Conversely, let us consider the subgroup consisting of all
elements of \[G(\hat\C)\] which fix two specified points \[z_0\ne z_1\].
It is convenient to conjugate by an automorphism
which carries \[z_0\] to zero and \[z_1\] to infinity. The argument above
shows that an automorphism \[g\] fixes both zero and infinity if and only
if it has the form \[g(z)=\lambda z\] for some \[\lambda\ne 0\]. Thus the
subgroup consisting of all such elements is isomorphic to the multiplicative
group \[\C\ssm\{0\}\], and hence is commutative as required.

Finally, consider automorphisms \[g\] which fix only the point at infinity.
A similar argument shows that \[g(z)\] must be equal to \[z+c\]
for some constant \[c\ne 0\]. (Compare 1.13.)
These automorphisms, together with the identity map, form a subgroup of
\[G(\hat\C)\] which is isomorphic to the additive group of \[\C\]; which again
is commutative as required. Further details will be left to the reader.\QED

The corresponding discussion for the group \[G(H)\] (or \[G(D)\])
will be important in \S3. To fix our ideas, let
us concentrate on the case of the half-plane. It follows from 1.7 that every
automorphism of \[H\] extends uniquely to an automorphism of \[\hat
\C\]. Hence 1.12 applies immediately. However, we must consider not only
fixed points inside of \[H\] but also fixed points which lie on the
boundary \[\partial H=\R\cup\infty\]. A priori, we should also consider
fixed points which lie in the lower half-plane, completely outside of the
closure \[\bar H\]. However, it is easy to check that \[w\] is a fixed point
of a fractional linear transformation with real coefficients
if and only if the complex
conjugate \[\bar w\] is also a fixed point. Thus each fixed point in the lower
half-plane, outside of \[\bar H\], is paired with a fixed point
which is inside the upper half-plane \[H\].\smallskip

{\bf 1.14. Definition.} The non-identity automorphisms of \[H\] fall into
three classes, as follows:

An automorphism \[g\in G(H)\] is said to be {\bit elliptic} if it has a
fixed point \[w_0\in H\]. We may also describe \[g\] as a {\bit rotation}
around \[w_0\]. (Compare (1) above.)

The automorphism \[g\in G(H)\] is {\bit hyperbolic} if it has two distinct
fixed points on the boundary \[\partial H\]. As an example,
\[g\] fixes the two points \[0\]
and \[\infty\] if and only if it has the form \[g(w)=\lambda w\] with
\[\lambda>0\].

The automorphism \[g\in G(H)\]
is {\bit parabolic} if it has just one double fixed
point, which necessarily belongs to the boundary \[\partial H=\R\cup
\infty\]. For example, if this double fixed point is the point at infinity,
then \[g\] is necessarily a translation: \[g(w)=w+c\;\] for some constant
\[c\ne 0\] in \[\R\].

{\QP{\bf 1.15. Lemma.} \it Two non-identity elements of \[G(H)\] commute if
and only if they have exactly
the same fixed points in \[\bar H=H\cup\R\cup\infty\]. The set
of all group elements with some specified
fixed point set, together with the identity
element, forms a commutative subgroup, which is isomorphic to a circle in the
elliptic case and to a real line in the parabolic or hyperbolic cases.
\medskip}

The proof is easily supplied.\QED

Again, we could equally well work with \[D\] in place of \[H\].
One defect of this exposition is that it requires going outside of the
half-plane \[H\] in order to distinguish between parabolic and hyperbolic
automorphisms. For a more intrinsic description of the difference between
these cases, see Problem 1-5.

\lln\medskip

We conclude this section with a number of problems for the reader.\medskip

{\bf Problem 1-1.} If a holomorphic map \[f:D\to D\] fixes the
origin, and is not a rotation, prove that the successive images \[f^
{\circ n}(z)\] converge to zero for all \[z\] in the open disk \[D\],
and prove that this convergence is uniform on compact subsets of \[D\].
(Here \[f^{\circ n}\] stands for the \[n$-fold iterate \[f\circ\cdots\circ
f\]. The example \[f(z)=z^2\] shows that convergence need not be uniform
on all of \[D\].) \smallskip

{\bf Problem 1-2.} Show that the group \[G(D)\] of conformal
automorphisms of the unit disk \[D\] consists of all maps
$$	g(z)\;=\; e^{i\theta} (z-a)/(1-\bar a z)	$$
with \[|e^{i\theta}|=1\], where \[a\in D\] is the unique point which maps
to zero. (Compare 1.7.) Check that \[|g'(0)|\le 1\], and conclude using 1.2
that any holomorphic map \[f:D\to D\] must satisfy \[|f'(0)|\le 1\].
\smallskip

{\bf Problem 1-3.} Show that the action of the group
\[G(\hat\C)\] on \[\hat\C\]
is {\bit simply 3-transitive}. That is, there is one and only one automorphism
which carries three distinct specified points of \[\hat\C\]
into three other specified
points. Similarly, show that the action of of \[G(\C)\] on \[\C\] is simply
2-transitive. Show that the action of \[G(D)\] on the boundary circle
\[\partial D\] carries three specified points into three other specified
points if and only if they have the same cyclic order; and show that the
action of \[G(D)\] carries two points of \[D\] into two other specified
points if and only if they have the same \Poincare distance.\smallskip

%{\bf Problem 1-4.} Using polar coordinates \[z=re^{i\theta}\] on the unit
%disk, show that the \Poincare metric is given by
%$$	ds=2|dz|/(1-r^2) =2\sqrt{(dr^2+r^2d\theta^2)}/(1-r^2)\ge 2|dr|
%/(1-r^2)\,,	$$
%where equality holds along any radial line \[\theta={\rm constant}\].
%Conclude that the associated distance function satisfies
%\[\PD(0,\, re^{i\theta})=\log((1+r)/(1-r))\]. In particular,
%show that this expression tends to infinity as \[r\to 1\], and hence check
%directly that the
%boundary of \[D\] has infinite \Poincare distance from any interior point, and
%that the closed neighborhood \[N_r(0,\rho_D)\] is a compact subset of \[D\].
%(\S1.8.)\smallskip

{\bf Problem 1-4.} By an {\bit anti-holomorphic} automorphism of \[\hat\C\]
we mean an orientation reversing self-homeomorphism of the form \[z\mapsto
\eta(\bar z)\], where \[\eta\] is holomorphic.
If \[L\] is a straight line or circle in
\[\hat\C\], show that there is one and only one anti-holomorphic involution
of \[\hat\C\] having \[L\] as fixed point set, and show that no other
non-vacuous fixed point sets can occur. Show that the automorphism group
\[G(\hat\C)\] acts transitively on the set of straight lines and circles
in \[\hat\C\]. Similarly, show that any
anti-holomorphic involution of \[D\] is the reflection in some \Poincare
geodesic; and show that \[G(D)\] acts transitively on the set of such
geodesics.\smallskip

{\bf Problem 1-5.}  Let \[g\] be an automorphism of \[D\] with \[g\circ g\]
not the identity map. Show that \[g\] is hyperbolic
if and only if there exists an automorphism \[h\] which satisfies
$$	h\circ g\circ h^{-1}\;=\;g^{-1}\,,	$$
or if and only if there exists some necessarily unique geodesic \[L\] with
respect to the \Poincare metric which is mapped onto itself by \[g\],
or if and only if \[g\] commutes with some anti-holomorphic involution.
(The possible choices for \[h\] are just the \[180
^\circ\] rotations about the points of \[L\].)\smallskip

{\bf Problem 1-6.} Define the {\bit infinite band} \[B\subset\C\] of width
\[\pi\] to be the set
of all \[z=x+iy\] with \[|y|<\pi/2\]. Show that the exponential map carries
\[B\] isomorphically onto the right half-plane \[\{u+iv:u>0\}\]. Show that
the \Poincare metric on \[B\] takes the form
$$	ds=|dz|/cos\,y\,. \eqno (7)	$$
Show that the
real axis is a geodesic whose \Poincare arclength coincides with its usual
Euclidean arclength, and show that each real translation \[z\mapsto z+c\] is
a hyperbolic automorphism of \[B\] having the real axis as its unique
invariant geodesic.\smallskip
\eject

{\bf Problem 1-7.} Given four distinct points \[z_j\] in
\[\hat\C\] show that the {\bit cross ratio}
$$	\chi(z_1\,,\,z_2\,;\,z_3\,,\,z_4)
	\;\;=\;\;{(z_1-z_3)(z_2-z_4)\over(z_1-z_4)(z_2-z_3)}\;\;\in
	\;\;\C\ssm\{0,1\} $$
is invariant under fractional linear transformations, and show that \[\chi\]
is real if and
only if the four points lie on a straight line or circle. \smallskip

{\bf Problem 1-8.} Show that each \Poincare
neighborhood \[N_r(w_0\,,\,\rho_H)\] in the upper half-plane is bounded
by a Euclidean circle, but that \[w_0\] is not its Euclidean center. Show that
each \Poincare geodesic in the upper half-plane is a straight line or
semi-circle which meets the real axis orthogonally. If the geodesic 
through \[w_1\] and \[w_2\] meets \[\partial H=\R\cup\infty\] at the points
\[\alpha\] and \[\beta\], show that the \Poincare distance between \[w_1\]
and \[w_2\] is equal to the logarithm of the cross ratio
\[\chi(w_1\,,\,w_2\,;\,\alpha\,,\,\beta)\]. Prove corresponding
statements for the unit disk.\smallskip

{\bf Problem 1-9.} The Gaussian curvature of a conformal metric \[ds=
\gamma(w)|dw|\] with \[w=u+iv\] is given by the formula
$$	 K\;\;=\;\; {\gamma_u^2+\gamma_v^2 -\gamma(\gamma_{uu}+\gamma_{vv})
	\over \gamma^4}	$$
where the subscripts stand for partial derivatives. (Compare Willmore, p. 79.)
Check that the \Poincare metrics (3), (4) and (7) have curvature \[K\equiv-1\],
and more generally that the metric
\[ds=c|dw|/v\] has curvature \[K\equiv-1/c^2\].
Check that the spherical metric (6) has curvature \[K\equiv +1\].\smallskip

{\bf Problem 1-10.} Classify conjugacy classes in the groups \[G(H)\cong
\PSL(2,\R)\] as follows.
Show that every automorphism of \[H\] without fixed point is conjugate to
a unique transformation of the form \[w\mapsto w+1\] or \[w\mapsto w-1\] or
\[w\mapsto\lambda w\] with \[\lambda>1\]; and show that the conjugacy class of
an automorphism \[g\] with fixed point \[w_0\in H\] is uniquely determined
by the derivative \[\lambda=g'(w_0)\], where \[|\lambda|=1\].
Show also that each non-identity element of \[\PSL_2(\R)\] belongs
to one and only one ``one-parameter subgroup'', and that each
one-parameter subgroup is conjugate to either
$$      t\;\;\mapsto\;\;
        \left[\matrix{1&t\cr 0&1\cr}\right]\quad{\rm or}\quad
        \left[\matrix{e^t&0\cr 0&e^{-t}\cr}\right]\quad{\rm or}\quad
        \left[\matrix{\cos t&\sin t\cr -\sin t&\cos t\cr}\right]      $$
according as its elements are parabolic or hyperbolic or elliptic.
Here \[t\] ranges over the additive group of real numbers.\smallskip

{\bf Problem 1-11.} For a non-identity automorphism \[g\in G(\hat\C)\], show
that the derivatives \[g'(z)\] at the two
fixed points are reciprocals, say \[\lambda\] and \[\lambda^{-1}\],
and show that the sum \[\lambda+\lambda^{-1}\] is a complete conjugacy
class invariant. (Here \[\lambda=1\] if and only if the two fixed points
coincide. In the
special case of a fixed point at infinity, one must evaluate the
derivative using the local coordinate \[\zeta=1/z\].)\smallskip

{\bf Problem 1-12.} Show that the conjugacy class of
a non-identity automorphism \[\;g(z)=\lambda z+c\;\]
in the group \[G(\C)\] is uniquely
determined by its image under the homomorphism \[g\mapsto \lambda\in\C\ssm
\{0\}\].

\vfil
%\rightline{jwm, version of 8-28-91}
\eject
% \bigskip\bigskip
\footline={\hss\tenrm 2-\folio\hss}\pageno=1

\centerline{\bf \S2. The  Universal Covering, Montel's Theorem.}\bigskip

If \[S\] is a completely arbitrary Riemann surface, then the universal
covering \[\tilde S\] is a well defined simply connected Riemann surface,
with a canonical projection map \[p:\tilde S\to S\].
(Compare Munkres, and also Appendix E.) According to the Uniformization
Theorem, this universal covering \[\tilde S\] must be
conformally isomorphic to one of the three model surfaces (\S1.1).
Thus we have the following.

{\QP{\bf 2.1. Lemma.} \it Every Riemann surface \[S\] is conformally
isomorphic to a quotient of the form \[\tilde S/\Gamma\], where
\[\tilde S\] is a simply connected surface which is isomorphic to
either \[D\], \[\C\], or \[\hat\C\]. Here \[\Gamma\] %subset G(\tilde S)\]
is a discrete group of conformal automorphisms of \[\tilde S\],
such that every non-identity element of \[\Gamma\] acts without fixed points
on \[\tilde S\].\medskip}

This discrete group \[\Gamma\subset G(\tilde S)\] can be identified with
the fundamental group \[\pi_1(S)\]. The elements of \[\Gamma\]
are called {\bit deck transformations}. They can be characterized as
maps \[\gamma:\tilde S\to\tilde S\] which satisfy the identity
\[\gamma=p\circ\gamma\], that is maps for which the diagram
$$      \matrix{\tilde S&\buildrel \gamma\over\longrightarrow& \tilde S\cr
       \;\;\;\downarrow p &&\;\;\;\downarrow p\cr
        S&\buildrel = \over\longrightarrow& S}  $$
is commutative.
Conversely, if we are given a group \[\Gamma\] of conformal automorphisms
of a simply connected surface \[\tilde S\] which is discrete as a subgroup of
\[G(\tilde S)\], and such that every non-identity element of \[\Gamma\] acts
without fixed points,
then it is not difficult to check that \[\Gamma\] is the group
of deck transformations for a universal covering map \[\tilde S\to\tilde
S/\Gamma\]. (Compare Problem 2-1.)\smallskip

We can analyze the three possibilities for \[\tilde S\] as follows.\smallskip
% (Note that (1) implies the sharper condition that \[\Gamma\] acts
% properly discontinuously --- that is, for any compact subset \[K\subset\tilde
% S\] there are only finitely many \[\gamma\] with \[\gamma(K)\cap K\ne
% \emptyset\].)

{\bf Spherical Case.} According to 1.12, every conformal automorphism of the
Riemann sphere \[\hat\C\] has
at least one fixed point. Therefore, if \[S\cong\hat\C/\Gamma\]
is a Riemann surface with universal covering
\[\tilde S\cong \hat\C\], then the group \[\Gamma\subset G(\hat\C)\] must be
trivial, hence \[S\] itself must be isomorphic to \[\hat\C\].\smallskip

{\bf Euclidean Case.} By 1.13, the group \[G(\C)\] of conformal
automorphisms of the complex plane
consists of all affine transformations \[z\mapsto \lambda z+c\] with
\[\lambda\ne 0\]. Every such transformation with \[\lambda
\ne 1\] has a fixed point.
Hence, if \[S\cong\C/\Gamma\] is a surface with universal covering
\[\tilde S\cong \C\], then \[\Gamma\] must be a discrete 
group of translations \[z\mapsto z+c\] of the complex plane \[\C\].
There are three subcases:

If \[\Gamma\] is trivial, then \[S\] itself is isomorphic to \[\C\].

If \[\Gamma\] has just one generator, then \[S\] is isomorphic to
the cylinder \[\C/\Z\], which in turn is isomorphic under the exponential
map to the {\bit punctured plane} \[\C\ssm\{0\}\].

If \[\Gamma\] has two generators, then it can be described as a two-dimensional
{\bit lattice}\break
\[\Lambda\subset\C\], that is an additive group generated by
two complex numbers which are linearly independent over \[\R\]. The
quotient \[{\bf T}=\C/ \Lambda\] is called a {\bit torus}.\smallskip

In all three of these cases, note that our surface inherits a locally
Euclidean geometry from the Euclidean metric \[|dz|\] on its universal
covering surface. For example the punctured plane \[\C\ssm\{0\}\], consisting
of points \[\exp(z)=w\], has a complete locally Euclidean metric \[|dz|
=|dw/w|\]. It is easy to check that such a metric is unique up to a
multiplicative constant. (Compare Problem 2-2.) We will refer to these
Riemann surfaces as [complete locally] {\bit
Euclidean surfaces}. The term ``{\bit parabolic surface\/}'' is also commonly used
in the literature.\smallskip

{\bf Hyperbolic Case.} In all other cases, the universal covering \[\tilde
S\] must be conformally isomorphic to the unit disk. Such Riemann surfaces
are said to be {\bit Hyperbolic}. {\it As an example, any Riemann surface
with non-abelian fundamental group, and in particular any
surface of higher genus, is necessarily Hyperbolic.}

{\bf Remark.} Here the word ``Hyperbolic'' is a reference to Hyperbolic
Geometry, that is the geometry of Lobachevsky.
Unfortunately the term ``hyperbolic'' has at least three
different well established meanings in conformal dynamics. (Compare \S1.14
and \S14.) In a crude attempt to avoid confusion, I will alway capitalize the
word when it is used with this geometric meaning.\smallskip

Every Hyperbolic surface \[S\] possesses a unique {\bit \Poincare metric},
which is complete, with Gaussian curvature identically equal to \[-1\].
To construct this metric, we note that the \Poincare metric on \[\tilde S\]
is invariant under the action of \[\Gamma\]. Hence there is one and only one
metric on \[S\] so that the projection \[\tilde S\cong D\to S\]
is a local isometry. Hence, just as in \S1.9, there is an
associated {\bit \Poincare distance function} \[\rho(z,z')=\rho_S(z,z')\] which
is equal to the length of a shortest geodesic from \[z\]to \[z'\]. As in 1.10
we have:

{\QP{\bf 2.2. Lemma.} \it If \[f:S\to T\] is a holomorphic map
between Hyperbolic surfaces, then
$$	 \PD_T(f(z),\,f(z'))\;\le\;\PD_S(z,z')\,.	$$
Furthermore, if equality holds for some \[z\ne z'\], then it follows
that \[f\] is a {\bit local isometry}. That is, \[f\] preserves the
infinitesimal distance element \[ds\], and hence preserves the distances
between nearby points.\medskip}

{\bf Proof.} This follows, just as in 1.10, since a minimal geodesic
joining \[z\] to \[z'\] must be mapped isometrically.\QED

{\bf Caution:} It no longer follows that \[f\] must be a conformal
isomorphism from \[S\] to \[T\]. However,
if \[f\] is a local isometry, then we can at least assert that \[f\] induces
an isometry \[\tilde S\buildrel\cong\over\longrightarrow\tilde T\] between
the universal covering surfaces, and hence that \[f\] is a covering map
from \[S\] onto \[T\].\smallskip

Here is an example. The map \[f(z)= z^2\] on the
{\bit punctured disk} \[D\ssm\{0\}\] is certainly not an automorphism, since
it is
two-to-one. However, the universal covering of \[D\ssm\{0\}\] can be identified
with the left half-plane, mapped to \[D\ssm\{0\}\] by the exponential map,
and \[f\] lifts to the automorphism \[w\mapsto 2w\] of this left half-plane.
Therefore, \[f\] is a covering map, and
preserves the \Poincare metric locally.\smallskip

Note that the punctured disk has abelian fundamental group \[\pi_1(D\ssm\{0\})
\cong\Z\]. Here is another such example.
For any \[ r>1\], the {\bit annulus}
$$	{\cal A}_r\;=\;\{\,z\; : \; 1<|z|<r\,\}		$$
is a Hyperbolic surface, since it admits a holomorphic map to the unit disk.
The fundamental group \[\pi_1({\cal A}_r)\] is evidently also free cyclic.
In fact annuli and the punctured disk
are the only Hyperbolic surfaces with abelian fundamental
group, other than the disk itself. (Compare Problem 2-3.)\smallskip

{\bf Maximal Hyperbolic Example.} If \[a_1\,,\,a_2\,,
\,a_3\] are three distinct points of \[\hat\C\], then the complement
\[\;\Sigma_3\,=\,\hat\C\ssm\{a_1\,,\,a_2\,,\,a_3\}\,\cong\,\C\ssm\{0,1\}\;\]
is called a {\bit thrice punctured sphere}. This is evidently a Hyperbolic
surface, for example since its fundamental group is not abelian.
One immediate corollary of this observation is the following.

{\QP {\bf 2.3. Picard's Theorem.}
{\it Any holomorphic map from \[\C\] to \[\hat\C\] which omits three
different values must be constant. More generally, if the Riemann surface
\[S\] admits some non-constant
holomorphic map to the thrice punctured
sphere \[\Sigma_3\], then \[S\] must be Hyperbolic.}\medskip}

For this map \[f:S\to \Sigma_3\] can be
lifted to a holomorphic map from the universal covering \[\tilde S\]
to the universal covering \[\tilde \Sigma_3\cong D\]. By Liouville's
Theorem (\S1.3) it follows that \[\tilde S\cong D\].\QED

Let \[U\subset \hat\C\] be any connected open set which omits at least
three points, and hence is Hyperbolic.
It is often useful to compare the \Poincare distance between two points
of \[U\] with the spherical distance between the same two points within
\[\hat\C\]. (See \S1(6).) Here is a crude estimate.
Again let \[N_r(z,\rho_U)\] be the closed neighborhood of
some fixed radius \[r\] with respect to the \Poincare metric about
the point \[z\] of \[U\].

{\QP {\bf 2.4. Lemma.} {\it 
As \[z\] converges towards the boundary \[\partial U\] in the spherical
metric,  the spherical diameter
of the neighborhood \[N_r(z,\rho_U)\] tends to zero.}\medskip}

\noindent (For a sharper statement in the simply connected case, see
A.8 in the Appendix.) \smallskip

{\bf Proof of 2.4.}
First consider the special case \[U=\Sigma_3=\hat\C\ssm\{a_1\,,\,
a_2\,,\,a_3\}\]. Choose some fixed base point \[z_0\] in \[U=\Sigma_3\].
For each fixed \[r>0\], the neighborhood \[N_r(z_0,\,\rho_U)\] is the image
under projection of a corresponding neighborhood in the universal covering
surface, and hence is compact and connected. (Compare 1.9.)
Now, as \[z\] tends to one of the three boundary points \[a_1\] of \[U\],
it must eventually
leave any compact subset of \[U\], hence the distance \[\PD_U(z,\,z_0)\] must
tend to infinity. For fixed \[r\], it follows that the neighborhood
\[N_r(z,\,\rho_U)\] will eventually be
disjoint from any given compact subset of \[\Sigma_3\]. In fact, since the set
\[N_r(z,\,\rho_U)\] is connected, this entire set must tend to just one
boundary point \[a_1\] with respect to the spherical metric.

Now consider an arbitrary Hyperbolic open set \[U\subset\hat\C\]. Given a
sequence of points of \[U\] tending to the boundary, we can choose a
subsequence \[\{z_j\}\] which converges to a single boundary point \[a_1\].
Choose two other boundary points \[a_2\]
and \[a_3\], and consider the inclusion map
from \[U\] to \[\Sigma_3=\hat\C\ssm\{a_1\,,\,a_2\,,\,a_3\}\]. Applying the
Pick inequality 2.2, we have \[N_r(z_j\,,\,\rho_U)\subset N_r(z_j\,,\,\rho_
{\Sigma_3})\], and it follows from the discussion above that
this entire neighborhood converges to the boundary point \[a_1\]
as \[j\to\infty\].\QED\smallskip

Using the \Poincare metric, we will develop another important tool.
A sequence of maps \[f_n:S\to\hat\C\] on a Riemann surface
\[S\] is said to {\bit converge locally uniformly} (or\break {\bit uniformly
on compact sets}) to the limit \[g:S\to\hat\C\] if for
every compact subset \[K\subset S\] the\eject
\noindent sequence
\[\{f_n|K\}\] of maps \[f_n\] restricted to \[K\] converges
uniformly to \[g|K\]. Here it is to be understood that we use the
spherical metric \[\SD(z_1\,,\,z_2)\] on the target space \[\hat\C\].
\smallskip

{\bf Definition.} Let \[S\] and \[S'\] be Riemann surfaces, with \[S'\]
compact. A collection \[{\cal F}\] of holomorphic maps
\[f_\alpha:S\to S'\] is {\bit normal}\[\] if every infinite
sequence of maps from \[{\cal F}\] contains a subsequence which converges
locally uniformly to a limit.\smallskip

(For the case of a non-compact target space \[S'\], see 2.6 below.)
Note that the limit function \[g\] must itself be holomorphic, by the Theorem
of Weierstrass. However, this limit \[g\] need not belong to the given family.
Roughly speaking, a family of maps is normal if and only if its closure,
in the space of all holomorphic maps from \[S\] to \[S'\], is a compact
set. (Ahlfors [A1, p.213].)

{\QP {\bf 2.5. Montel's Theorem.} {\it Let \[S\] be any Riemann surface,
which we may as well suppose to be Hyperbolic. If a collection \[{\cal F}\]
of holomorphic maps from \[S\] to the Riemann sphere \[\hat\C\] takes values
in some Hyperbolic open subset \[U\subset\hat\C\], or equivalently
if there are three distinct points of \[\hat\C\] which never occur as
values, then this collection \[{\cal F}\] is normal.} \medskip}

More explicitly, any sequence of holomorphic maps \[f_n:S\to U\]
contains a subsequence which converges, uniformly on compact subsets of \[S\],
to some holomorphic map \[g:S\to\bar U\]. Here it is essential that \[g\] be
allowed to take values in the closure \[\bar U\]. However, the proof will
show that the image \[g(S)\] is either contained in \[ U\], or is a single
point belonging to the boundary of \[ U\].\smallskip

{\bf Proof of 2.5.} First note by
2.3 that the surface \[S\] must be Hyperbolic, unless all of our maps
are constant. Hence \[S\] also has a \Poincare metric.
Choose a countable dense subset of points \[z_j\in S\,,\;j\ge 1\].
(It follows easily from \S1.1 that every Riemann surface possesses a countable
dense subset. This statement was first proved by Rad\'o. Compare
Ahlfors \& Sario.) Given any sequence
of holomorphic maps \[f_n:S\to U\subset\hat\C\], we can first choose an
infinite subsequence \[\{f_{n(p)}\}\] of the \[f_n\] so that the
images \[f_{n(p)}(z_1)\] converge to a limit within the closure of
\[U\]. Then choose a sub-sub-sequence
\[f_{n(p(q))}\] so that the \[f_{n(p(q))}(z_2)\] also converge to a limit,
and continue inductively. By a diagonal procedure,
taking the first element of the first subsequence, the second element of the
second subsequence, and so on, we construct a new infinite sequence of maps
\[g_m=f_{n_m}\] so that \[\lim_{m\to\infty}g_m(z_j)\in \bar U\] exists
for every choice of \[z_j\].

{\bf Case 1.} Suppose that every one of these limit points in \[\bar U\]
actually belongs to the set \[U\] itself. Given any compact set \[K\subset S\]
and any \[\epsilon>0\], we can choose finitely many \[z_j\] so that every
point \[z\in K\] has Poincar\'e distance \[\rho_S(z,z_j)<\epsilon\] from one of
these \[z_j\]. Further, we can choose \[m_0\] so that 
\[\rho_U(g_m(z_j)\,,\,g_n(z_j))<\epsilon\] for each of these finitely many
\[z_j\], whenever \[m,\,n>m_0\,\]. For any \[z\in K\] it then follows using
2.2 that \[\PD_U(g_m(z)\,,\,g_n(z))<3\epsilon\] whenever \[m,\,n>m_0\].
Thus the \[g_m(z)\] form a Cauchy sequence.
It follows that the sequence of functions \[\{g_m\}\] converges
to a limit, and that this convergence is uniform
on compact subsets of \[S\].

{\bf Case 2.} Suppose on the other hand that \[\lim_{m\to\infty}
g_m(z_j)\] is actually a boundary point \[a_0\in\partial U\] for some \[z_j\].
Then it follows
from 2.4 that \[g_m(z)\] converges to \[a_0\] for every \[z\in S\],
and that this convergence is uniform on compact subsets of \[S\].\QED
 \smallskip

{\bf 2.6. Remark.} If we consider maps \[S\to S'\] where the target space
\[S'\] is not compact, then the definition should be modified as follows.
We continue to assume that \[S\] is connected.
A collection \[\cal F\] of maps \[f_\alpha:S\to S'\] is {\bit normal~}
if every infinite sequence of maps from \[\cal F\] contains a subsequence
which either

\noindent $(1)$  converges locally uniformly
to a holomorphic map from \[S\] to \[S'\], or

\noindent $(2)$ diverges locally uniformly to infinity, in the sense that
the successive images of any compact subset of \[S\] eventually miss any given
compact subset of \[S'\].\smallskip

\lln

{\bf Concluding Problems.}\smallskip

{\bf Problem 2-1.} Let \[S\] be a simply connected Riemann surface, and let
\[\Gamma\subset G(S)\] be a discrete subgroup; that is, suppose that the
identity element
is an isolated point of \[\Gamma\]. If every non-identity element of
\[\Gamma\] acts on \[S\] without fixed points, show that the action of
\[\Gamma\] is
{\bit properly discontinuous}. That is, for every compact \[K\subset S\]
show that
only finitely many group elements satisfy \[K\cap\gamma(K)\ne
\emptyset\]. Show that each \[z\in S\] has a neighborhood \[U\] whose
translates \[\gamma(U)\] are pairwise disjoint. Conclude that \[S/\Gamma\]
is a well defined Riemann surface with \[S\] as its universal covering.
\smallskip

{\bf Problem 2-2.} Show that any Riemann surface which
is not conformally isomorphic to \[\hat\C\],
has one and up to a multiplicative constant only
one conformal Riemann metric which is complete, with constant Gaussian
curvature. (Make use of Hopf's Theorem which asserts that, for each
real number \[K\], there is one and only
one complete simply connected surface of constant curvature \[K\] up to
isometry. See Willmore p. 162.) On the other hand, show that \[\hat\C\]
has a 3-dimensional family of conformal metrics with curvature \[+1\].
\smallskip

{\bf Problem 2-3.} Using Problems 1-10 and 1-6, show that every Hyperbolic
surface
with abelian fundamental group is conformally isomorphic either to the disk
\[D\], or the punctured disk \[D\ssm\{0\}\], or to the annulus
$$	{\cal A}_r\;=\;\{z\in\C:1<|z|<r\}	$$
for some uniquely defined \[ r>1\]. Define the {\bit modulus} of
this annulus to be the number \[\mod\!({\cal A}_r)
=\log r/2\pi>0\].
Show that this annulus has a unique simple closed geodesic, which has
length \[\ell=\pi/\mod\!({\cal A}_r)=2\pi^2/\log r\,\]. On the other hand,
show that the punctured disk \[D\ssm\{0\}\] has no closed geodesic.
This punctured disk, which is conformally isomorphic to the set
\[{\cal A}_\infty\;=\;\{z\in\C:1<|z|<\infty\}\], is sometimes described
as an annulus of infinite modulus. (However this designation is ambiguous,
since the Euclidean surface \[\C-\{0\}\]
might also be described as an annulus of infinite modulus.)
\smallskip %By definition, this punctured disk has infinite modulus.\smallskip

{\bf Problem 2-4.} If \[S\] and \[S'\] are Hyperbolic Riemann surfaces
(not necessarily compact), show that {\bit every~} family of maps from
\[S\] to \[S'\] is normal.\smallskip

{\bf Problem 2-5.}
Show that normality is a local property. More precisely,
let \[S\] and \[S'\] be any Riemann surfaces,
and let \[\{f_\alpha\}\] be a family of holomorphic maps from \[S\] to \[S'\].
If every point of \[S\] has a neighborhood \[U\] such
that the collection \[\{f_\alpha|U\}\] of restricted maps
is normal, show by a diagonal argument as in the
proof of 2.5 that \[\{f_\alpha\}\] is normal.
\bigskip\vfill

%\rightline{jwm, version of 8-28-91}
\end
